\documentclass{tum_report}
\usepackage[T1]{fontenc}
\usepackage{combelow}
\usepackage{booktabs}
\usepackage{tabularx}
\usepackage{microtype}
\counterwithout{table}{section}
\counterwithin{table}{chapter}
\hypersetup{pdftitle={Laboratory Work No. 3: Romanian Playfair Cipher},
    pdfauthor={Mihail Pancenco},pdfsubject={Cryptography, Task 3.1, variant 19}}
\lstdefinestyle{reportcode}{
    language=[Sharp]C,basicstyle=\ttfamily\fontsize{9}{11}\selectfont,
    keywordstyle=\bfseries\color{blue!45!black},
    commentstyle=\itshape\color{black!55},stringstyle=\color{green!35!black},
    columns=fullflexible,keepspaces=true,showstringspaces=false,
    breaklines=true,breakatwhitespace=false,emptylines=0,numbers=none,
    frame=single,rulecolor=\color{black!25},backgroundcolor=\color{black!2},
    aboveskip=6pt,belowskip=6pt,
    literate={Ă}{{\u{A}}}1 {ă}{{\u{a}}}1 {Â}{{\^{A}}}1 {â}{{\^{a}}}1
      {Î}{{\^{I}}}1 {î}{{\^{i}}}1 {Ș}{{\cb{S}}}1 {ș}{{\cb{s}}}1
      {Ț}{{\cb{T}}}1 {ț}{{\cb{t}}}1 {Ş}{{\c{S}}}1 {ş}{{\c{s}}}1
      {Ţ}{{\c{T}}}1 {ţ}{{\c{t}}}1
}
\lstdefinestyle{console}{style=reportcode,language={},emptylines=1,
    basicstyle=\ttfamily\fontsize{9}{11}\selectfont,breakatwhitespace=true}
\newcommand{\codepart}[2]{{\setstretch{1}\lstinputlisting[
    style=reportcode,firstline=#1,lastline=#2,gobble=4]{../code/playfair/PlayfairCipher.cs}}}
\newcommand{\consolepart}[1]{{\setstretch{1}\lstinputlisting[
    style=console]{evidence/#1.txt}}}
\clubpenalty=10000
\widowpenalty=10000

\begin{document}
\ministry{MINISTRY OF EDUCATION AND RESEARCH\\OF THE REPUBLIC OF MOLDOVA}
\university{\textbf{Technical University of Moldova}}
\faculty{\textbf{Faculty of Computers, Informatics and Microelectronics}}
\department{\textbf{Department of Software Engineering and Automatics}}
\title{Cryptography}{}
\reportyear{2026}

\begin{titlepage}
\centering
\includegraphics[width=0.20\textwidth]{UTM_Logo.png}\par\vspace{0.8cm}
{\bfseries
\textsc{\ministryname}\\[0.15cm]
\universityname\\
\facultyname\\
\departmentname\par}

\vspace{1cm}
\textbf{Mihail Pancenco}\quad\textbf{FAF-243}\par
\vspace{0.9cm}
{\Large\textbf{Report}\par}
\vspace{0.5cm}
{\Large\normalfont\textit{Laboratory Work No. 3}\par}
\vspace{0.35cm}
{\Large\normalfont\textit{Cryptography}\par}
\vspace{0.7cm}
{\huge\normalfont\textit{Playfair Cipher}\par}
\vspace{0.4cm}
{\large\normalfont Romanian alphabet\\Task 3.1 -- Variant 19\par}

\vspace{1cm}
\begin{flushright}
\normalfont
Supervisors:\par
\textbf{Poștaru Andrei}, \textit{Doctor}\par
\textbf{Maia Zaica}, \textit{University assistant}\par
DISA, FCIM, UTM
\end{flushright}
\vfill
\textbf{Chișinău -- \degreeyear}\par
\end{titlepage}

\clearpage
\addtocontents{toc}{\protect\thispagestyle{empty}}
\tableofcontents
\thispagestyle{empty}
\clearpage
\pagestyle{fancy}

\chapter{Task and short theory}
\section{Task 3.1 and the Romanian alphabet}
The objective of Laboratory Work No. 3 is to study classical ciphers that
use more than a fixed substitution of individual letters. My number in the
group list is 19. Since it is odd, the selected task is 3.1: implement the
Playfair cipher for Romanian messages, as specified in the assignment~\cite{ref:assignment}.
The program must accept encryption and decryption, require a key of at
least seven letters, reject invalid input clearly, and leave the final
restoration of word spaces to the user.

Playfair encrypts pairs of letters using a matrix built from a keyword.
The Romanian alphabet has 31 letters. In this implementation, J is treated
as I, so the matrix has 30 cells arranged in \textbf{six rows and five columns}.
The letter Î remains distinct from I. Uppercase and lowercase Romanian
letters are accepted; the cedilla variants of Ș and Ț are normalized.

The keyword's first occurrences are written into the matrix first. All
unused letters are then appended in Romanian alphabetical order. Messages
are converted to uppercase, ordinary spaces are removed, and J is replaced
by I. Identical letters within a pair are separated with X; Q is used
when the repeated letter is X. An unmatched final letter receives the
same choice of filler. Decryption keeps these fillers for manual interpretation.

\section{Encryption and decryption rules}
Let $A[r,c]$ be the matrix, using zero-based coordinates
$r\in\{0,\ldots,5\}$ and $c\in\{0,\ldots,4\}$. Let $d=1$ for encryption
and $d=-1$ for decryption. For a pair at positions $(r_1,c_1)$ and $(r_2,c_2)$:
\begin{itemize}
\item \textbf{Same row:} keep each row and replace each column by
      $(c+d+5)\bmod5$.
\item \textbf{Same column:} keep each column and replace each row by
      $(r+d+6)\bmod6$.
\item \textbf{Rectangle:} use the opposite corners $(r_1,c_2)$ and $(r_2,c_1)$.
      This rule is identical for encryption and decryption.
\end{itemize}
Thus encryption moves right or down, while decryption moves left or up.
The modulo operations implement wraparound. Adding the relevant dimension
keeps the remainder nonnegative during decryption.

Playfair substitutes letter pairs, which obscures simple single-letter
frequency relationships. It is nevertheless a classical educational cipher,
not a suitable method for protecting modern confidential communication.

\clearpage
\chapter{Implementation}
\section{Program structure, normalization and validation}
The C\# implementation contains two classes. \texttt{Program} manages the
menu, reads input, catches validation errors and repeats the corresponding
prompt. \texttt{PlayfairCipher} builds the matrix, prepares letter pairs,
encrypts and decrypts. The project targets .NET 10. The excerpts below
are taken directly from the current source~\cite{ref:source}; empty lines
are compressed for readability.

\codepart{100}{125}
\noindent\textit{Comment.} \texttt{Normalize} replaces cedilla variants,
converts text to uppercase, merges J into I and removes ordinary spaces.
It then checks every remaining character against the 30 matrix symbols.
Digits, punctuation, tabs and other unsupported characters are rejected.
The error identifies the rejected character and states the accepted letters.

\codepart{15}{25}
\noindent\textit{Comment.} The constructor forbids spaces in the keyword,
then normalizes it and enforces a minimum length of seven letters. This
check occurs \emph{before} duplicate removal; repeated letters are allowed
in a valid keyword. Empty messages are rejected by \texttt{Encrypt}, and
\texttt{Decrypt} requires a nonempty, even-length ciphertext after normalization.

\clearpage
\section{Building the keyword matrix}
\codepart{5}{9}
\noindent\textit{Comment.} The matrix alphabet contains 30 symbols because
J uses the I cell. The dimensions are six rows and five columns.

\codepart{127}{159}
\noindent\textit{Comment.} The first loop retains each keyword letter once,
in order of first appearance. The second loop appends unused alphabet
letters. The nested loops fill the matrix from left to right, row by row.
This produces exactly one occurrence of every matrix symbol.

For example, the key \texttt{criptografie} becomes
\texttt{CRIPTOGRAFIE}; its distinct letters are \texttt{CRIPTOGAFE}.
The resulting matrix is
\begin{center}\setstretch{1.15}
\begin{tabular}{ccccc}\toprule
C & R & I/J & P & T\\
O & G & A & F & E\\
Ă & Â & B & D & H\\
Î & K & L & M & N\\
Q & S & Ș & Ț & U\\
V & W & X & Y & Z\\\bottomrule
\end{tabular}\end{center}
\texttt{FindPosition} locates a letter by scanning the rows and columns.
\texttt{PrintMatrix} displays the same cells, labelling the I cell as I/J.

\clearpage
\section{Preparing pairs and inserting fillers}
\texttt{PreparePairs} scans the normalized message using an index. The
first letter of the next pair is \texttt{message[index]}. The filler is
X, except when that letter is itself X, in which case the filler is Q.
The following branches decide the second letter and how far to advance:
\codepart{172}{190}
\noindent\textit{Comment.} If only one letter remains, the method appends
a filler and consumes one original letter. If the next two letters are
identical, it forms a pair with a filler and also advances by only one:
the second original letter must be reconsidered. Otherwise, it consumes
both letters. Each constructed pair is added to a \texttt{List<string>}.

\begin{table}[H]\caption{Pair preparation for BALLOON}
\centering
\begin{tabular}{cllc}\toprule
Index & Next letters & Pair produced & Next index\\\midrule
0 & BA & BA & 2\\
2 & LL & LX & 3\\
3 & LO & LO & 5\\
5 & ON & ON & 7\\\bottomrule
\end{tabular}\end{table}
The prepared text is \texttt{BALXLOON}, or \texttt{BA LX LO ON}.
Only repeated letters inside the current pair are separated. For example,
\texttt{ABBA} already splits into \texttt{AB BA} and needs no filler.
For \texttt{XX}, the method produces \texttt{XQ XQ}; for a final X, it
produces \texttt{XQ}. This prevents an identical-letter filler pair.

\clearpage
\section{Transforming pairs and reversing the cipher}
\codepart{196}{225}
\noindent\textit{Comment.} \texttt{FindPosition} supplies the coordinates
of the two input letters. The three branches implement the same-row,
same-column and rectangle rules. \texttt{Encrypt} passes direction 1;
\texttt{Decrypt} passes direction $-1$. The result letters are read from
the matrix, and one short console line shows the input pair, output pair
and rule used.

During encryption, the message is normalized and sent through
\texttt{PreparePairs}. During decryption, the ciphertext is normalized and
split directly into two-letter substrings; no fillers are inserted.
Both methods concatenate the transformed pairs and return a string.

The console prints the normalized key, distinct keyword letters, matrix,
normalized message and prepared pairs before the transformations. Filler
insertions are reported where they occur. The final output is either the
ciphertext or the recovered prepared plaintext. Word spaces, capitalization,
the original I/J choice and ambiguous filler letters are not reconstructed
automatically.

\clearpage
\chapter{Results and verification}
\section{Manual check of all three rules}
The project builds successfully with zero compiler warnings and errors.
The console transcripts in this chapter were captured from the running,
unchanged program. For a matrix that is especially easy to verify by hand,
I used the seven-letter key \texttt{AĂÂBCDE}. Its matrix follows the base
alphabet order:
\begin{center}\setstretch{1.15}
\begin{tabular}{c|ccccc}
 & 1 & 2 & 3 & 4 & 5\\\hline
1 & A & Ă & Â & B & C\\
2 & D & E & F & G & H\\
3 & I/J & Î & K & L & M\\
4 & N & O & P & Q & R\\
5 & S & Ș & T & Ț & U\\
6 & V & W & X & Y & Z\\
\end{tabular}\end{center}
The coordinates in Table~\ref{tab:rules} start at one for convenient manual
checking. The code uses zero-based array indices.
\begin{table}[H]\caption{Manual encryption and inverse checks}
\label{tab:rules}\centering\small
\begin{tabular}{cllcc}\toprule
Input & Input coordinates & Rule & Cipher pair & Decrypts to\\\midrule
AC & (1,1), (1,5) & Same row & ĂA & AC\\
AV & (1,1), (6,1) & Same column & DA & AV\\
AȚ & (1,1), (5,4) & Rectangle & BS & AȚ\\\bottomrule
\end{tabular}\end{table}
For AC, C wraps from column 5 to column 1. For AV, V wraps from row 6 to
row 1. For AȚ, the opposite rectangle corners contain B and S.

The diagnostic input \texttt{AC AV AȚ} exercises all three rules in one run:
\consolepart{rules_encrypt}
Decrypting \texttt{ĂA DA BS} with the same key gives \texttt{ACAVAȚ}.
This confirms the inverse left/up shifts and the unchanged rectangle rule.

\clearpage
\section{Repeated letters and decryption}
For key \texttt{AĂÂBCDE} and input \texttt{bal loon}, the program removes
the input space, converts the text to uppercase and inserts X between
the repeated L letters:
\consolepart{balloon_encrypt}

With the same key, decrypting \texttt{CĂ KY ÎQ PO} returns the prepared
plaintext rather than deleting the inserted X:
\consolepart{balloon_decrypt}
The result \texttt{BALXLOON} agrees with the prepared encryption input.
The user can identify the inserted X and restore the intended word spacing.
Removing every X automatically would also remove genuine plaintext letters.

For the boundary case \texttt{XX}, the program prints
\texttt{Pairs: XQ XQ} and produces \texttt{YPYP}. For a single input A,
it appends X, prints \texttt{Pairs: AX} and produces \texttt{ÂV}.
These results confirm both filler selection and final-letter padding.

\clearpage
\section{Romanian text and a permuted keyword matrix}
I also used the keyword \texttt{criptografie} and the Romanian input
\texttt{ştiinţă şi ţară}. The cedilla variants are accepted and normalized
to comma-below letters; the resulting matrix is the one shown in Section 2.2.
\consolepart{romanian_encrypt}
The ciphertext is \texttt{UIAITLQDXAȘFCÂ}. Decrypting it with the same key
returns \texttt{ȘTIXINȚĂȘIȚARĂ}: the same normalized message with the
inserted X retained. This run checks Romanian diacritics, lowercase input,
spaces, duplicate keyword letters and a keyword-dependent alphabet order.

\section{Validation and input recovery}
The program rejects invalid input and repeats the corresponding prompt.
Table~\ref{tab:validation} summarizes captured validation runs; the full
console transcripts are saved in \texttt{evidence/}.
\begin{table}[H]\caption{Input validation cases and observed responses}
\label{tab:validation}\centering\small
\begin{tabularx}{\textwidth}{l l X}\toprule
Field & Input & Response\\\midrule
Keyword & \texttt{secret} & Minimum length is seven letters before removing duplicates.\\
Keyword & \texttt{cripto1} & Character 1 is rejected; the permitted Romanian letters are stated.\\
Keyword & \texttt{cripto grafie} & Spaces in the keyword are rejected.\\
Message & Empty input & At least one Romanian letter is required.\\
Message & \texttt{salut!} & Character ! is rejected; the permitted Romanian letters are stated.\\
Message & \texttt{salut2} & Character 2 is rejected; the permitted Romanian letters are stated.\\
Ciphertext & \texttt{ABC} & A nonzero, even number of letters is required.\\\bottomrule
\end{tabularx}\end{table}
After invalid keys, supplying \texttt{AĂÂBCDE} continues the operation.
After invalid message characters, \texttt{Salut} is accepted. After the
odd-length ciphertext \texttt{ABC}, supplying \texttt{ĂADABS} successfully
decrypts to \texttt{ACAVAȚ}. Thus an error does not terminate the program.

\clearpage
\chapter{Conclusions}
In this laboratory work, I implemented the Romanian Playfair cipher using
a keyword matrix with six rows and five columns. Merging J into I made
room for the 31-letter Romanian alphabet while keeping Î as a separate
letter. The program handles uppercase and lowercase input, Romanian
cedilla variants and ordinary message spaces.

I verified the three transformation rules and their inverses: horizontal
shifts, vertical shifts and rectangle substitution. The examples also
confirmed the importance of preparing pairs correctly. When two current
letters are identical, only the first original letter is consumed; the
second must be reconsidered. X is the usual filler, with Q used when an
X would otherwise create an identical pair.

The console output makes the process easy to check: it shows the matrix,
prepared pairs and one short explanation for each transformation.
Decryption correctly recovers the prepared plaintext. Restoring word
spaces, distinguishing an original J from I and resolving filler letters
remain manual tasks because this information is not preserved by the
cipher. The work reinforced my understanding of pair substitution, modular
wraparound and input validation through a small, readable implementation.

\section{Verification coverage and source code}
For each of two keyed matrices, I checked all $30\times29=870$ ordered
pairs of distinct letters against an independent coordinate calculation,
then decrypted the results. Both sets passed: 1,740 pair cases in total.
The checks also covered repeated letters, alternate fillers, odd-length
plaintext, Romanian normalization, the minimum key length, invalid input
and decryption back to the prepared plaintext. The verification run passed
3,537 assertions. The source files were not modified while preparing this report.

Source files: \texttt{Program.cs} and \texttt{PlayfairCipher.cs},
in \texttt{lab\_03/code/playfair/}. The project repository is
\begin{center}
\url{https://github.com/MiHCraft/Cryptography_UTM}
\end{center}

\vspace{0.6cm}
\phantomsection
\addcontentsline{toc}{chapter}{Bibliography}
\begin{thebibliography}{9}
\singlespace\normalsize\raggedright
\bibitem{ref:assignment}
Technical University of Moldova.
\textit{Lucrare de laborator nr. 3: Cifruri polialfabetice}.
Course assignment, Section 3.2 (Playfair), pp.~3--6, and Task 3.1, p.~7.

\bibitem{ref:source}
Pancenco, Mihail.
\textit{Romanian Playfair Cipher}. C\# implementation,
\texttt{lab\_03/code/playfair/Program.cs} and
\texttt{lab\_03/code/playfair/PlayfairCipher.cs}.
Project repository: \url{https://github.com/MiHCraft/Cryptography_UTM}.
\end{thebibliography}

\end{document}
